\subsection{\texorpdfstring{\(\mathbf{T}^n\) 的自同构}{T 的 n 次幂的自同构}}

设 H 是 \(R^{n}\) 的闭子群， \(\varphi\) 是 \(R^{n}\) 到商群 \(R^{n}/H\) 上的典范同态。若 f 是 \(R^{n}/H\) 到拓扑群 G 中的连续同态，则 \(\dot{f} = f \circ \varphi\) 是 \(R^{n}\) 到 G 中的连续同态，它是周期的，且其周期群包含 H；反过来，每个周期群包含 H 的 \(R^{n}\) 到 G 中的周期连续同态都具有这种形式。

在 \(H = Z^{n}\) 的情形，商群 \(R^{n}/Z^{n} = T^{n}\) 是紧的，因此只要 G 是豪斯多夫的， \(T^{n}\) 到拓扑群 G 中的每个连续同态 f 都是 \(T^{n}\) 到 G 中的严格态射（第 III 章 § 2 第 8 小节注 1）； \(\dot{f} = f \circ \varphi\) 是 \(R^{n}\) 到 G 中的严格态射；此外， \(f(\mathbf{T}^{n}) = \dot{f}(\mathbf{R}^{n})\) 是 G 的一个紧子群，同构于某个群 \(T^{p}\) （ \(o \leqslant p \leqslant n\) ）。

特别地，我们看到 \(T^{n}\) 到群 \(R^{m}\) 中的连续同态只有零映射，因为 \(\{0\}\) 是 \(R^{m}\) 的唯一紧子群。

让我们把这个应用于 \(\mathbf{T}^n\) 到群 \(\mathbf{T}^p\) 中的连续同态；若 \(f\) 是这样一个同态， \(\varphi\) 是 \(\mathbf{R}^n\) 到 \(\mathbf{T}^n\) 上的典范同态，则 \(f \circ \varphi\) 是 \(\mathbf{R}^n\) 到 \(\mathbf{T}^p\) 中的连续同态；因而（第 3 小节命题 3），若用 \(\psi\) 表示 \(\mathbf{R}^p\) 到 \(\mathbf{T}^p\) 上的典范同态，则存在线性映射 \(u\) （ \(\mathbf{R}^n\) 到 \(\mathbf{R}^p\) 中的），使得 \(f \circ \varphi = \psi \circ u\) 。若 \(\boldsymbol{x} \in \mathbf{Z}^n\) ，则 \(f(\varphi(\boldsymbol{x}))\) 是 \(\mathbf{T}^p\) 的单位元，故必有 \(u(\boldsymbol{x}) \in \mathbf{Z}^p\) ，即 \(u(\mathbf{Z}^n) \subset \mathbf{Z}^p\) 。反过来，若 \(u\) 是满足此条件的 \(\mathbf{R}^n\) 到 \(\mathbf{R}^p\) 中的任意线性映射，则 \(\psi \circ u\) 是 \(R^n\) 到 \(T^p\) 中的周期连续同态，其周期群包含 \(Z^n\) ，因此定义了 \(T^n\) 到 \(T^p\) 中的一个连续同态。

考虑在什么条件下 \(f\) 是 \(\mathbf{T}^n\) 到 \(\mathbf{T}^p\) 的一个子群上的同构。首先， \(u\) 必须是 \(\mathbf{R}^n\) 到 \(\mathbf{R}^p\) 中的单射映射，否则向量子空间 \(\overline{u}^{-1}(\mathrm{o})\) 就会包含任意接近原点的点 \(x \neq 0\) ，而在这样的点处我们将有

\[
f (\varphi (x)) = f (\varphi (0)) \quad \text { and } \quad \varphi (x) \neq \varphi (0),
\]

与假设矛盾。这一条件蕴含 \(p \geqslant n\) 。这时像 \(u(\mathbf{Z}^n)\) 是一个秩为 \(n\) 的离散子群（群 \(\mathbf{Z}^p\) 的）； \(u(\mathbf{Z}^n)\) 关于 \(\mathbf{Z}^p\) 的不变因子（§ 1 第 1 小节）必须全等于 1，否则就会存在一点 \(x \in \mathbf{Z}^n\) 和一个整数 \(k > 1\) ，使得 \(u(k^{-1}\mathbf{x}) \in \mathbf{Z}^p\) 而 \(k^{-1}\mathbf{x} \notin \mathbf{Z}^n\) ，从而 \(f(\varphi(k^{-1}\mathbf{x})) = f(\varphi(0))\) 且 \(\varphi(k^{-1}\mathbf{x}) \neq \varphi(0)\) ，与假设矛盾。反之，若这一条件满足，则 \(u(\mathbf{R}^n) \cap \mathbf{Z}^n = u(\mathbf{Z}^n)\) ，且 \(f\) 是 \(\mathbf{T}^n\) 到 \(u(\mathbf{R}^n)/u(\mathbf{Z}^n)\) 上的一个同构。

若把这一论证应用于 p = n 的情形，我们就得到下列命题：

命题 5. 拓扑群 \(\mathbf{T}^n\) 到其一个子群上的每个同构都是 \(\mathbf{T}^n\) 的一个自同构，这个自同构是由一个线性映射 \(u\) （ \(\mathbf{R}^n\) 到它自身上的）过渡到商而得到的，该映射限制在 \(\mathbf{Z}^n\) 上时是 \(\mathbf{Z}^n\) 的一个自同构。

等价地（§ 1 第 1 小节），若 \(u(\boldsymbol{e}_i) = \sum_{j=1}^{n} a_{ij} \boldsymbol{e}_j\) ，则诸 \(a_{ij}\) 必须是使 \(\det(a_{ij}) = \pm 1\) 成立的整数。特别地，对于 \(n = 1\) ：

命题 6. 拓扑群 T 到其一个子群上的同构只有恒等映射和对称 \(x \to -x\) 。

